\section{Trees}

So far, we've only looked at trees on $\omega$ or trees on some other set. These are subtrees of $\parenth{{^{<\omega}\Omega, \subsetneq}}$. Here is the general notion.

\begin{boxdefinition}[Tree]
    $\parenth{T, <_T}$ is a \textbf{tree} iff it's a strict partial ordering (transitive and irreflexive) and, for all $y \in T$, $\setst{x}{x <_T y}$ is well-ordered by $<_T$.
\end{boxdefinition}

Fix a tree $\parenth{T, <_T}$.

Each node of a tree sits at some height, and the nodes at the same height make up a level.

\begin{boxdefinition}[Height and Levels]
    For $y \in T$, define
    \begin{align*}
        \height_T(y) &= \type\of{\setst{x}{x <_T y}, <_T} \\
        \level_{\alpha}(T) &= \setst{y \in T}{\height_T(y) = \alpha} \\
        \height(T) &= \text{the least $\alpha$ such that $\level_{\alpha}(T) = \emptyset$}
    \end{align*}
    We call $\height_T(y)$ the \textbf{height} of $y$, $\level_{\alpha}(T)$ the \textbf{$\alpha$th level} of $T$, and $\height(T)$ the \textbf{height} of $T$.
\end{boxdefinition}

We will mostly care about trees that are as well-behaved as possible.

\begin{boxdefinition}[Normal Tree]
    We say $\parenth{T, <_T}$ is \textbf{normal} iff
    \begin{enumerate}
        \item $T$ has a unique root
        \item Splitting to the next level: given $x$ with $\height(x) = \alpha$, there are $y, z \in \level_{\alpha + 1}(T)$ such that $y \neq z$ and $x <_T y, z$.
        \begin{figure}[H]
            \centering
            \begin{tikzpicture}[dot/.style = {circle, fill, inner sep = 1.5pt}, every node/.style = {font = \small}]
                \draw[dashed, gray] (-2, 0) -- (2, 0) node[right, black] {$\level_{\alpha}(T)$};
                \draw[dashed, gray] (-2, 1.5) -- (2, 1.5) node[right, black] {$\level_{\alpha + 1}(T)$};
                \draw[dotted] (0, -1) -- (0, 0);
                \node[dot, label = {right:$x$}] (x) at (0, 0) {};
                \node[dot, label = {above:$y$}] (y) at (-0.8, 1.5) {};
                \node[dot, label = {above:$z$}] (z) at (0.8, 1.5) {};
                \draw (x) -- (y);
                \draw (x) -- (z);
            \end{tikzpicture}
        \end{figure}
        \item Extensions at every higher level: $\forall \alpha < \height(T)$, $\forall \beta$ with $\alpha < \beta < \height(T)$, $\forall x \in \level_{\alpha}(T)$, $\exists y \in \level_{\beta}(T)$ such that $x <_T y$.
        \item Extensionality at limit levels\footnote{so far, in homework, we've only seen trees with at most $\omega$ many levels; here, we're not assuming this of every tree}: given a limit ordinal $\alpha < \height(T)$, if $y, z \in \level_{\alpha}(T)$ and $\setst{x}{x <_T y} = \setst{x}{x <_T z}$, then $y = z$. So the following picture, in which $y \neq z$ sit on a limit level above exactly the same nodes, cannot occur in a normal tree.
        \begin{figure}[H]
            \centering
            \begin{tikzpicture}[dot/.style = {circle, fill, inner sep = 1.5pt}, every node/.style = {font = \small}]
                \foreach \h in {0, 1, 1.7, 2.15, 2.45}
                    \node[dot] at (0, \h) {};
                \draw (0, 0) -- (0, 2.45);
                \node at (0, 2.85) {$\vdots$};
                \draw[dashed, gray] (-2, 4) -- (2, 4) node[right, black] {$\level_{\alpha}(T)$, $\alpha$ limit};
                \node[dot, label = {above:$y$}] (y) at (-0.8, 4) {};
                \node[dot, label = {above:$z$}] (z) at (0.8, 4) {};
                \draw[dotted] (0, 3.2) -- (y);
                \draw[dotted] (0, 3.2) -- (z);
                \node[left, align = right] at (-0.3, 1.1) {$\setst{x}{x <_T y}$ \\ $= \setst{x}{x <_T z}$};
            \end{tikzpicture}
        \end{figure}
    \end{enumerate}
\end{boxdefinition}

Finally, branches are the paths through a tree, of any length.

\begin{boxdefinition}[Branch]
    A \textbf{branch} of a tree $\parenth{T, <_T}$ is a subset $b \subset T$ such that $\forall y \in b$, $\forall x <_T y$, $x \in b$ and $b$ is linearly ordered (equivalently, well-ordered) by $<_T$. We say $b$ is an \textbf{$\alpha$-branch} if $\type\of{b, <_T} = \alpha$.
\end{boxdefinition}

Two nodes of a tree may or may not lie on a common branch.

\begin{boxdefinition}[Comparability and Incomparability]
    Given a tree $T$, we say $x$ and $y$ are \textbf{comparable} (written $x \parallel y$) iff $x \leq_T y$ or $y \leq_T x$. Else, we say $x$ and $y$ are \textbf{incomparable} and write $x \perp y$.
\end{boxdefinition}

Sets of pairwise incomparable nodes will matter a great deal.

\begin{boxdefinition}[Antichain]
    $A \ssq T$ is an \textbf{antichain} (often abbreviated ac) if $\forall x, y \in A$, $x \neq y$ implies $x \perp y$.
\end{boxdefinition}

The levels are the basic example.

\begin{boxexample}
    For every $\alpha$, $\level_{\alpha}\of{T}$ is an antichain.
\end{boxexample}

An antichain may also be as large as possible.

\begin{boxdefinition}[Maximal Antichain]
    An antichain $A \ssq T$ in a tree $T$ is \textbf{maximal} iff
    \begin{enumerate}
        \item $A$ is an antichain
        \item There's no antichain $B$ with $B \supsetneq A$
    \end{enumerate}
\end{boxdefinition}

Note: if $A$ is a maximal antichain then for every $x \in T$, there's some $a \in A$ such that $a \parallel x$. Otherwise, $A \cup \set{x} = B$ would be an antichain.
